% VSEPR-SIM VERSION REPORT
% State Vector Audit, Core Implementation Review, and Beta-6 Pre-Repair Manifest
% Version: v4.0.4.05 (audit) / v5.0.0-beta6 (migration target)
% Generated: 2026

\documentclass[11pt,a4paper]{article}
\usepackage{amsmath, amssymb, geometry, booktabs, hyperref, xcolor, enumitem, titlesec, fancyhdr, listings, mdframed}
\geometry{margin=2.5cm}
\hypersetup{colorlinks=true, linkcolor=blue, urlcolor=blue}
\pagestyle{fancy}
\fancyhf{}
\rhead{VSEPR-SIM Version Report}
\lhead{v4.0.4.05 $\to$ v5.0.0-beta6}
\cfoot{\thepage}

\definecolor{red1}{HTML}{CC0000}
\definecolor{green1}{HTML}{006600}
\definecolor{amber}{HTML}{B36200}
\definecolor{gray1}{HTML}{444444}
\definecolor{codebg}{HTML}{F5F5F5}

\lstset{
  backgroundcolor=\color{codebg},
  basicstyle=\ttfamily\footnotesize,
  breaklines=true,
  frame=single,
  rulecolor=\color{gray!50}
}

\title{
  \textbf{VSEPR-SIM Version Report}\\[0.5em]
  \large State Vector Audit \textemdash{} Core Implementation Review\\
  \large F/U Check Analysis \textemdash{} Update Loop Survey \textemdash{} Beta-6 Repair Manifest\\[1em]
  \normalsize \texttt{v4.0.4.05} $\longrightarrow$ \texttt{v5.0.0-beta6}
}
\author{VSEPR-SIM Development \textemdash{} Internal Technical Document}
\date{2026 \textemdash{} Generated pre-refactor}

\begin{document}
\maketitle
\tableofcontents
\newpage

% ==============================================================================
\section{Document Purpose}
% ==============================================================================

This document is the mandatory pre-refactor audit for the \texttt{v5.0.0-beta6} migration.
It describes:

\begin{enumerate}
  \item The core-of-core implementation: what the engine actually does right now.
  \item The full state vectors for both the atomistic and coarse-grained bead layers.
  \item Force (\textbf{F}) and energy (\textbf{U}) check mechanisms as implemented.
  \item The update loop structure and ensemble types.
  \item The macro-solving methods available.
  \item All detected defects and inconsistencies from the audit.
  \item The immediate action table with priority classification.
\end{enumerate}

This document is \textbf{not aspirational}. It describes what exists, not what is planned.
Planned changes are listed only in Section~\ref{sec:action_table}.

% ==============================================================================
\section{Core-of-Core: What the Engine Actually Does}
% ==============================================================================

\subsection{The Fundamental Architecture}

VSEPR-SIM is structured as a layered state-transformation engine:

\[
S_{n+1} = \mathcal{K}(S_n,\, \Delta t,\, \Lambda)
\]

where $\mathcal{K}$ is the kernel operator, $\Delta t$ is the timestep, and
$\Lambda \in \{\text{C-Empirical}, \text{C-Reactive}, \text{C-Polarizable},
\text{CG-Reduced}, \text{CG-Enriched}, \text{Hybrid-AdRes}\}$ is the model fidelity level.

The layers are:

\begin{center}
\begin{tabular}{lll}
\toprule
Layer & Namespace & Role \\
\midrule
Primitive math   & \texttt{vsepr::}       & Vec3, geom\_ops, flat coord arrays \\
Atomistic state  & \texttt{atomistic::}   & State, StateC, integrators, models \\
Energy kernel    & \texttt{vsepr::pot::}  & Bond, angle, torsion, LJ, VSEPR domain \\
Optimizer        & \texttt{vsepr::}       & FIRE (geometry minimisation) \\
Coarse-grain     & \texttt{coarse\_grain::} & Bead, BeadSystem, mapping \\
V5 transport     & \texttt{v5::}          & BeadDynState, transport shell \\
Thermodynamics   & \texttt{vsepr::thermo::} & Gibbs, ThermoDatabase \\
\bottomrule
\end{tabular}
\end{center}

\subsection{Two Parallel Engine Tracks}

There are \textbf{two distinct coordinate conventions} operating simultaneously.

\textbf{Track A: \texttt{src/sim} + \texttt{src/pot}} \\
Coordinates as flat \texttt{std::vector<double>}, stride 3.
Used by: FIRE optimizer, bond/angle/torsion/LJ/VSEPR energy kernels.
This is the VSEPR geometry engine.

\textbf{Track B: \texttt{atomistic/}} \\
Coordinates as \texttt{std::vector<Vec3>}.
Used by: velocity Verlet, Langevin, FIRE (atomistic variant), LJ+Coulomb MD model.
This is the molecular dynamics engine.

Both tracks define their own \texttt{Vec3}. They do not interoperate.
This is the primary structural defect of the current codebase.

% ==============================================================================
\section{State Vectors}
% ==============================================================================

\subsection{Atomistic Runtime State: \texttt{atomistic::State}}

The canonical per-step simulation state. Lives in \texttt{atomistic/core/state.hpp}.

\[
S = (N,\; \mathbf{X},\; \mathbf{V},\; \mathbf{T},\; \mathbf{Q},\; \mathbf{M},\;
     \mathrm{type},\; B,\; L,\; \mathbf{F},\; E,\; \mathrm{box},\;
     \boldsymbol{\mu}^{\mathrm{ind}},\; \boldsymbol{\alpha},\; \boldsymbol{\mu}^0)
\]

\begin{center}
\begin{tabular}{llll}
\toprule
Symbol & Field & Type & Units \\
\midrule
$N$           & \texttt{N}                  & \texttt{uint32\_t}      & --- \\
$\mathbf{X}$  & \texttt{X[N]}               & \texttt{Vec3[]}         & \AA \\
$\mathbf{V}$  & \texttt{V[N]}               & \texttt{Vec3[]}         & \AA/fs \\
$\mathbf{T}$  & \texttt{T[N]}               & \texttt{double[]}       & K (proxy, optional) \\
$\mathbf{Q}$  & \texttt{Q[N]}               & \texttt{double[]}       & $e$ \\
$\mathbf{M}$  & \texttt{M[N]}               & \texttt{double[]}       & amu \\
              & \texttt{type[N]}            & \texttt{uint32\_t[]}    & species ID \\
$B$           & \texttt{B}                  & \texttt{Edge[]}         & bond graph \\
$L$           & \texttt{L}                  & \texttt{Event[]}        & event log \\
$\mathbf{F}$  & \texttt{F[N]}               & \texttt{Vec3[]}         & kcal/mol/\AA\; (scratch) \\
$E$           & \texttt{E}                  & \texttt{EnergyTerms}    & kcal/mol \\
              & \texttt{box}                & \texttt{BoxPBC}         & \AA \\
$\boldsymbol{\mu}^{\mathrm{ind}}$ & \texttt{induced\_dipoles[N]}   & \texttt{Vec3[]} & $e\cdot$\AA \\
$\boldsymbol{\alpha}$ & \texttt{polarizabilities[N]} & \texttt{double[]}  & \AA$^3$ \\
$\boldsymbol{\mu}^0$  & \texttt{permanent\_dipoles[N]}& \texttt{Vec3[]}  & $e\cdot$\AA \\
\bottomrule
\end{tabular}
\end{center}

\noindent\textbf{EnergyTerms ledger:}
\[
E_{\mathrm{total}} = U_{\mathrm{bond}} + U_{\mathrm{angle}} + U_{\mathrm{tors}}
                   + U_{\mathrm{vdW}} + U_{\mathrm{Coul}} + U_{\mathrm{ext}} + U_{\mathrm{pol}}
\]

All seven terms are individually tracked. The \texttt{total()} method sums them.
$U_{\mathrm{pol}}$ is the SCF polarization energy, wired to the induced dipole solver.

\noindent\textbf{BoxPBC:} Orthogonal minimum-image convention.
$\Delta\mathbf{r}_{ij} \leftarrow \Delta\mathbf{r}_{ij} - \mathbf{L} \cdot \mathrm{round}(\Delta\mathbf{r}_{ij}/\mathbf{L})$,
component-wise. Cached inverse $\mathbf{L}^{-1}$ avoids repeated division.
Triclinic cells: not supported.

\subsection{Physics-Informed Summary State: \texttt{atomistic::StateC}}

Formal notation (Dev Day $\sim$42):
\[
S_0 = \{m_0,\, E_0,\, \mathbf{x}_0,\, \mathbf{v}_0,\, C,\, \mathcal{E},\, K,\, \Sigma\}
\]
\[
S_f = \{m_f,\, E_f,\, \Phi,\, \mathcal{S},\, \Pi,\, \Omega,\, Q\}
\]
\[
S_f = T(S_0,\, t,\, \Lambda)
\]

\begin{center}
\begin{tabular}{lll}
\toprule
Symbol & Struct & Content \\
\midrule
$C$        & \texttt{Composition}         & Atomistic / Bead / Hybrid / Macro \\
$\mathcal{E}$ & \texttt{Environment}      & $T$, $P$, $\mathbf{E}$-field, $\mathbf{B}$-field, radiation, shear, flow \\
$K$        & \texttt{Constraints}         & Frozen atoms, bond break/form, symmetry \\
$\Sigma$   & \texttt{SourceTag}           & Origin, engine version, SHA-512 hash, parent ID \\
$\Lambda$  & \texttt{ModelLevel}          & 6-tier fidelity ladder \\
$\Phi$     & \texttt{StructuralState}     & Geometry class, topology, symmetry score \\
$\mathcal{S}$ & \texttt{StabilityMetric}  & $E/N$, $F_{\mathrm{RMS}}$, $F_{\mathrm{max}}$, local-minimum flag \\
$\Pi$      & \texttt{PerformanceProperties} & $D$, $\eta$, $\kappa$, $K_{\mathrm{bulk}}$, $\rho$, $T_m$ \\
$\Omega$   & \texttt{EventLog}            & Step-tagged events \\
$Q$        & \texttt{QualityFlags}        & Conservation booleans, drift values \\
\bottomrule
\end{tabular}
\end{center}

\subsection{Coarse-Grained Bead State: \texttt{coarse\_grain::Bead}}

\begin{center}
\begin{tabular}{lll}
\toprule
Field & Type & Role \\
\midrule
\texttt{position}          & \texttt{Vec3}                    & Bead center (\AA) \\
\texttt{velocity}          & \texttt{Vec3}                    & Bead velocity (\AA/fs) \\
\texttt{mass}              & \texttt{double}                  & $\sum m_i^{\mathrm{parent}}$ (amu) \\
\texttt{charge}            & \texttt{double}                  & $\sum q_i^{\mathrm{parent}}$ ($e$) \\
\texttt{type\_id}          & \texttt{uint32\_t}               & Index into BeadType table \\
$\Sigma_i$ & \texttt{StructuralRole}          & Inert / Ionic / Covalent / Metallic / Mixed \\
$\Lambda_i$ & \texttt{StabilityClass}         & Transient / Metastable / AmbientStable / BulkLattice \\
\texttt{parent\_atom\_indices} & \texttt{uint32\_t[]}         & Traceback into \texttt{atomistic::State} \\
\texttt{mapping\_rule\_id} & \texttt{uint32\_t}               & Producing mapping rule \\
\texttt{com\_position}     & \texttt{Vec3}                    & Center-of-mass reference \\
\texttt{cog\_position}     & \texttt{Vec3}                    & Center-of-geometry reference \\
\texttt{mapping\_residual} & \texttt{double}                  & $|\mathrm{COM} - \mathrm{COG}|$ (\AA) \\
\texttt{surface}           & \texttt{optional<SurfaceDescriptor>}   & Anisotropic shape \\
\texttt{multi\_channel}    & \texttt{optional<MultiChannelDescriptor>} & Enriched descriptor \\
\texttt{unified}           & \texttt{optional<UnifiedDescriptor>}  & Adaptive-resolution descriptor \\
\texttt{orientation}       & \texttt{BeadOrientation}         & Inertia-frame primary axis \\
\texttt{has\_orientation}  & \texttt{bool}                    & Orientation validity flag \\
\bottomrule
\end{tabular}
\end{center}

\subsection{V5 Transport Bead State: \texttt{v5::BeadDynState}}

Dynamic state for the V5 transport testbed.
Equation of motion:
\[
\frac{d\mathbf{v}}{dt} = A(\mathbf{u} - \mathbf{v}) + B\,\mathbf{G} + C\,\mathbf{H}
\]

\begin{center}
\begin{tabular}{lll}
\toprule
Field & Type & Content \\
\midrule
\texttt{position}, \texttt{velocity}, \texttt{acceleration} & \texttt{Vec3} & Full kinematic state \\
\texttt{n\_hat} & \texttt{Vec3} & Orientation axis (static, Phase A/B) \\
\texttt{env} & \texttt{EnvironmentState} & $(\eta,\, \rho,\, C,\, P_2)$ \\
\texttt{mass} & \texttt{double} & Reduced mass (default 1.0) \\
\texttt{a\_drag}, \texttt{a\_body}, \texttt{a\_grad}, \texttt{a\_total} & \texttt{Vec3} & Decomposed acceleration (diagnostics) \\
\texttt{mod\_report} & \texttt{ModulationReport} & Kernel modulation diagnostics \\
\bottomrule
\end{tabular}
\end{center}

% ==============================================================================
\section{F and U Checks}
% ==============================================================================

\subsection{Force Checks (F)}

\subsubsection{Analytical Gradient Verification}
Available in \texttt{src/sim/optimizer.hpp} via \texttt{check\_gradients = true}.
Central finite difference:
\[
\nabla_i E \approx \frac{E(\mathbf{x} + h\hat{e}_i) - E(\mathbf{x} - h\hat{e}_i)}{2h}
\]
$h = 10^{-6}$ \AA, tolerance $10^{-5}$ kcal/mol/\AA.
\textbf{Known bug:} the print path on max-error reporting re-calls \texttt{evaluate\_energy(coords)} instead of printing the stored numeric gradient, so the numeric column prints total energy. Diagnostics only --- non-fatal, but misleading.

\subsubsection{NaN / Inf Guard}
FIRE optimizer checks every iteration:
\begin{lstlisting}
if (has_invalid_values(result.coords) ||
    has_invalid_values(forces) ||
    !std::isfinite(result.energy))
    result.termination_reason = "NaN/Inf detected";
\end{lstlisting}
Terminates cleanly. Does not repair.

\subsubsection{Force Convergence Criteria}
Two simultaneous thresholds in \texttt{OptimizerSettings}:
\[
F_{\mathrm{RMS}} < 10^{-4}\;\frac{\mathrm{kcal}}{\mathrm{mol}\cdot\text{\AA}}, \quad
F_{\mathrm{max}} < 10^{-3}\;\frac{\mathrm{kcal}}{\mathrm{mol}\cdot\text{\AA}}
\]
Both must be satisfied simultaneously. These are tight relative to many production codes.

\subsubsection{Displacement Clamping}
Per-atom displacement clamped to \texttt{max\_step} = 0.2 \AA\ per step (vector magnitude, not component-wise). Prevents runaway steps on stiff geometries.

\subsubsection{FIRE Convergence: \texttt{atomistic} Track}
In \texttt{atomistic/integrators/fire.hpp}:
\[
F_{\mathrm{RMS}} < \epsilon_F = 10^{-6}, \quad \Delta U/N < \epsilon_U = 10^{-10}
\]
These are significantly tighter than the \texttt{src/sim} FIRE defaults and operate on different coordinate storage (Vec3 arrays vs flat doubles).

\subsection{Energy Checks (U)}

\subsubsection{EnergyBalance Ledger (\texttt{StateC})}
Enforces:
\[
E_0 + E_{\mathrm{in}} - E_{\mathrm{out}} \approx E_f + E_{\mathrm{loss}}, \quad
m_0 + m_{\mathrm{in}} - m_{\mathrm{out}} \approx m_f
\]
Violation sets \texttt{Q.energy\_conserved = false} and increments \texttt{Q.warning\_count}.
Does \textbf{not} halt simulation --- violation is logged, not fatal.

\subsubsection{EnergyTerms Completeness}
All seven terms are present in the ledger. \texttt{total()} sums all seven.
The VSEPR-mode EnergyResult struct also carries the full breakdown: bond, angle, torsion, nonbonded, VSEPR domain, vdW, Coulomb.

\subsubsection{Conservation in Velocity Verlet}
NVE conserves $E_{\mathrm{total}} = KE + PE$ by construction (symplectic integrator).
Tracked in \texttt{VelocityVerletStats}: \texttt{E\_drift = E\_final - E\_initial}.
No per-step assertion --- drift accumulates silently unless the caller inspects stats.

\subsubsection{Bead Conservation Checking}
\texttt{BeadSystem::check\_conservation()} compares $\sum m_b$ vs $\sum m_i^{\mathrm{atom}}$
and $\sum q_b$ vs $\sum q_i^{\mathrm{atom}}$ with tolerance $10^{-10}$.

% ==============================================================================
\section{Update Loop Structure}
% ==============================================================================

\subsection{Geometry Optimisation (FIRE, \texttt{src} track)}

Used for VSEPR structure generation. Not a dynamics integrator --- minimises $E(\mathbf{x})$.

\begin{enumerate}[label=\arabic*.]
  \item Compute $E$, $\mathbf{F} = -\nabla E$ at current $\mathbf{x}$.
  \item Check convergence ($F_{\mathrm{RMS}}$, $F_{\mathrm{max}}$, NaN/Inf, $dt_{\min}$).
  \item Compute power $P = \mathbf{F} \cdot \mathbf{v}$.
  \item FIRE velocity mixing: $\mathbf{v} \leftarrow (1-\alpha)\mathbf{v} + \alpha|\mathbf{v}|\hat{\mathbf{F}}$.
  \item Adaptive $dt$, $\alpha$: if $P > 0$, grow $dt$, decay $\alpha$; if $P \leq 0$, reset $\mathbf{v} = 0$, shrink $dt$, reset $\alpha$.
  \item Velocity Verlet half-step: $\mathbf{v}(t + \tfrac{\Delta t}{2}) = \mathbf{v}(t) + \mathbf{F}(t)\tfrac{\Delta t}{2}$.
  \item Position update: $\mathbf{x}(t+\Delta t) = \mathbf{x}(t) + \mathbf{v}(t+\tfrac{\Delta t}{2})\Delta t$. Clamp displacement.
  \item Evaluate $E$, $\mathbf{F}$ at new $\mathbf{x}$.
  \item Velocity Verlet second half: $\mathbf{v}(t+\Delta t) = \mathbf{v}(t+\tfrac{\Delta t}{2}) + \mathbf{F}(t+\Delta t)\tfrac{\Delta t}{2}$.
  \item Return to step 2.
\end{enumerate}

\subsection{NVE Molecular Dynamics (Velocity Verlet, \texttt{atomistic} track)}

Symplectic, time-reversible, conserves $E_{\mathrm{total}}$ (microcanonical ensemble).

\begin{enumerate}[label=\arabic*.]
  \item Half-kick: $\mathbf{v}_i \leftarrow \mathbf{v}_i + \dfrac{\mathbf{F}_i}{m_i}\dfrac{\Delta t}{2}$.
  \item Drift: $\mathbf{x}_i \leftarrow \mathbf{x}_i + \mathbf{v}_i \Delta t$.
  \item Force evaluation: \texttt{model.eval(state, params)} $\Rightarrow$ fills \texttt{F[N]}, \texttt{E}.
  \item Half-kick: $\mathbf{v}_i \leftarrow \mathbf{v}_i + \dfrac{\mathbf{F}_i}{m_i}\dfrac{\Delta t}{2}$.
  \item Diagnostics: $KE = \tfrac{1}{2}\sum m_i |\mathbf{v}_i|^2$, $T = \tfrac{2KE}{3Nk_B}$, log.
\end{enumerate}

\subsection{NVT Molecular Dynamics (Langevin, \texttt{atomistic} track)}

Stochastic equation of motion (Euler-Maruyama discretisation):
\[
m_i\,\frac{d\mathbf{v}_i}{dt} = \mathbf{F}_i - \gamma m_i \mathbf{v}_i
  + \sqrt{2\gamma m_i k_B T}\;\mathbf{R}(t)
\]
\begin{enumerate}[label=\arabic*.]
  \item Evaluate forces: \texttt{model.eval(state, params)}.
  \item Update velocity: $\mathbf{v}_i \leftarrow \mathbf{v}_i + [\mathbf{F}_i/m_i - \gamma\mathbf{v}_i]\Delta t
        + \sqrt{2\gamma k_B T/m_i}\sqrt{\Delta t}\,\mathbf{R}$.
  \item Update position: $\mathbf{x}_i \leftarrow \mathbf{x}_i + \mathbf{v}_i \Delta t$.
  \item Accumulate temperature statistics.
\end{enumerate}
Not symplectic --- energy is not conserved; temperature is controlled.

\subsection{V5 Transport Loop (\texttt{v5::BeadDynState})}

Mandatory update order (ERB §):
\begin{enumerate}[label=\arabic*.]
  \item Integrate poses $(\mathbf{R},\, \hat{\mathbf{n}})$ via velocity Verlet.
  \item Compute fast observables: $\rho$, $C$, $P_2$.
  \item Update slow state $\eta$.
  \item Evaluate modulated kernels $\mathcal{K}_k(\eta)$.
  \item Compute forces and torques.
  \item Return to step 1.
\end{enumerate}
Torque evolution and orientation dynamics are deferred (Phase A/B: rigid orientation).

\subsection{StateC Transformation Loop (\texttt{state\_c.hpp})}

\textbf{Status: stub.} Framework is complete; integrator is not connected.
Current loop body:
\begin{lstlisting}
// [PLACEHOLDER] In the real engine, force evaluation, integrator,
// thermostat, PBC, and constraint enforcement happen here.
\end{lstlisting}
The energy balance ledger, quality flags, event log, and provenance
chain are wired and functional. The force and position updates are not.

% ==============================================================================
\section{Simulation Types and Macro-Solving Methods}
% ==============================================================================

\subsection{Available Simulation Types}

\begin{center}
\begin{tabular}{lllll}
\toprule
Type & Ensemble & Integrator & Status & Track \\
\midrule
VSEPR geometry opt. & --- (minimisation) & FIRE & \textcolor{green1}{Implemented} & \texttt{src/sim} \\
NVE MD              & Microcanonical     & Velocity Verlet & \textcolor{green1}{Implemented} & \texttt{atomistic} \\
NVT MD              & Canonical          & Langevin        & \textcolor{green1}{Implemented} & \texttt{atomistic} \\
Geometry opt. (atomistic) & ---          & FIRE            & \textcolor{green1}{Implemented} & \texttt{atomistic} \\
V5 transport        & Non-equilibrium    & Velocity Verlet & \textcolor{amber}{Partial}      & \texttt{v5} \\
StateC evolution    & Configurable       & Stub            & \textcolor{red1}{Not connected} & \texttt{atomistic} \\
CG bead dynamics    & ---                & Not defined     & \textcolor{red1}{Not implemented} & \texttt{coarse\_grain} \\
\bottomrule
\end{tabular}
\end{center}

\subsection{Macro-Solving Methods}

\subsubsection{Gibbs Free Energy}
\[
G(T) = H_f - T \cdot S / 1000
\]
($H_f$ in kcal/mol, $S$ in cal/mol$\cdot$K.)
Seeded NIST reference database: 15 molecules, standard state 298.15 K / 1 atm.
Temperature correction is constant-$C_p$ approximation ($\int C_p\,dT$ not integrated).

\subsubsection{Equilibrium Constant}
\[
K_{\mathrm{eq}} = \exp\!\left(-\frac{\Delta G}{RT}\right)
\]

\subsubsection{SCF Polarization}
Iterative induced dipole solver:
\[
\boldsymbol{\mu}_i = \alpha_i \left(\mathbf{E}^0_i + \sum_{j \neq i} \mathbf{T}_{ij}\boldsymbol{\mu}_j \right)
\]
Wired into \texttt{State.Upol} and \texttt{EnergyTerms}.

\subsubsection{V4 Formation Engine}
Statistical mechanics of metallic formation.
Polynomial fitting ($11$--$15$ layer), eigen counters.
Continual formation from candidate central atoms (Au, Ag, Ni, Al, Fe, Cr, Ti, Co).
Version: \texttt{V4.0.4.04} (current), \texttt{V4.0.4.05} (audit target).

\subsubsection{Gas-Phase Thermodynamics (\texttt{gas2}, \texttt{gas3})}
Equation of state, kinetic theory, heat, nuclear corrections.
Sweep engine, polynomial fitting, quality scoring, state records.

% ==============================================================================
\section{Detected Errors and Inconsistencies}
\label{sec:errors}
% ==============================================================================

The following defects were identified during the pre-beta6 audit.

\subsection{E-1 \textemdash{} Two Vec3 Types (BLOCKER)}
\begin{mdframed}[backgroundcolor=red!8]
\textbf{Location:} \texttt{src/core/math\_vec3.hpp} (\texttt{vsepr::Vec3}) vs
\texttt{atomistic/core/state.hpp} (\texttt{atomistic::Vec3}).

\textbf{Problem:} Both types represent $\mathbb{R}^3$ with identical semantics but
are not interconvertible. Every interface between the energy kernel (Track A)
and the atomistic integrators (Track B) requires either a copy, a cast, or a
parallel reimplementation. This is not a latent bug --- it is an active barrier
to every planned cross-track integration.

\textbf{Impact:}
\begin{itemize}
  \item Coarse-grain beads import \texttt{atomistic::Vec3} for position/velocity.
  \item VSEPR energy kernels use \texttt{vsepr::Vec3} exclusively.
  \item No conversion path exists.
  \item Any function taking a bead position and feeding it to an energy kernel
        requires a full struct copy at the call site.
\end{itemize}

\textbf{Classification:} Blocker. Fix before refactor.
\end{mdframed}

\subsection{E-2 \textemdash{} \texttt{evolve()} Placeholder (BLOCKER)}
\begin{mdframed}[backgroundcolor=red!8]
\textbf{Location:} \texttt{atomistic/core/state\_c.hpp}, \texttt{evolve()} function.

\textbf{Problem:} The canonical state-transformation operator $T(S_0, t, \Lambda)$
contains a \texttt{[PLACEHOLDER]} comment where the integrator coupling belongs.
The loop runs for \texttt{n\_steps} but does not move particles or evaluate forces.
The energy balance always reports zero drift because the energy never changes.

\textbf{What the stub does correctly:}
\begin{itemize}
  \item Validates $S_0$ (mass $> 0$, energy finite).
  \item Initialises the energy balance ledger.
  \item Chains provenance on each step.
  \item Returns a structurally valid \texttt{OutcomeState}.
\end{itemize}

\textbf{What it does not do:}
\begin{itemize}
  \item Evaluate forces.
  \item Update positions or velocities.
  \item Advance the simulation by $\Delta t$.
\end{itemize}

\textbf{Classification:} Blocker. Core transformation is fiction without this.
\end{mdframed}

\subsection{E-3 \textemdash{} \texttt{sane()} Ignores \texttt{F[N]} Size}
\begin{mdframed}[backgroundcolor=amber!8]
\textbf{Location:} \texttt{atomistic/core/state.hpp}, \texttt{sane()} function.

\textbf{Problem:}
\begin{lstlisting}
inline bool sane(const State& s) {
    if (s.N == 0) return false;
    if (s.X.size() != s.N || s.V.size() != s.N ||
        s.Q.size() != s.N || s.M.size() != s.N ||
        s.type.size() != s.N) return false;
    return true;
}
\end{lstlisting}
\texttt{F[N]} is not checked. The force array can be empty or wrong size
and \texttt{sane()} returns \texttt{true}. \texttt{model.eval()} writes into
\texttt{F} --- if the array is undersized, this is undefined behaviour.

\textbf{Classification:} Medium. Cheap fix, prevents silent UB.
\end{mdframed}

\subsection{E-4 \textemdash{} \texttt{T[N]} Has No Update Rule}
\begin{mdframed}[backgroundcolor=amber!8]
\textbf{Location:} \texttt{atomistic/core/state.hpp}, field \texttt{T[N]}.

\textbf{Problem:} The per-particle temperature proxy \texttt{T[N]} is declared
in \texttt{State} but nothing in the visible codebase populates or updates it.
The correct system temperature is computed in \texttt{velocity\_verlet.hpp}
via the equipartition theorem:
\[
T = \frac{2KE}{3Nk_B}
\]
but is not written back into \texttt{T[N]}.

\textbf{Classification:} Medium. Mark as diagnostic-only until a
thermostat or local-temperature operator is implemented.
\end{mdframed}

\subsection{E-5 \textemdash{} Bead Mass Not Bridged to Transport Layer}
\begin{mdframed}[backgroundcolor=amber!8]
\textbf{Location:} \texttt{v5/bead\_transport.hpp}, \texttt{BeadDynState::mass}.

\textbf{Problem:} \texttt{BeadDynState.mass = 1.0} (reduced units, default).
Physical bead mass is:
\[
M_b = \sum_{i \in P_b} m_i
\]
stored correctly in \texttt{coarse\_grain::Bead::mass} (amu).
No conversion bridge exists between the physical mass and the reduced mass
used in the transport testbed. The equation of motion $\mathbf{a} = \mathbf{F}/m$
in the transport layer uses a mass that is not physically connected to the
source atomistic data.

\textbf{Classification:} Medium. Reduced units are acceptable; the
conversion must be explicit and stored.
\end{mdframed}

\subsection{E-6 \textemdash{} \texttt{StateC::from\_state()} Loses Metadata}
\begin{mdframed}[backgroundcolor=amber!8]
\textbf{Location:} \texttt{atomistic/core/state\_c.hpp}, \texttt{StateC::from\_state()}.

\textbf{Problem:} The factory correctly computes $m_0$, $E_0$, COM position,
COM velocity, and sets \texttt{comp = Atomistic}. It does not transfer:
\begin{itemize}
  \item \texttt{env} (Environment: temperature, pressure, fields, flow)
  \item \texttt{constraints} (Constraints: frozen atoms, bond rules, symmetry)
  \item \texttt{provenance.engine\_version}, \texttt{.hash}
\end{itemize}
The compressed state silently loses physical context. Any downstream
use of \texttt{StateC.env.temperature} after \texttt{from\_state()} returns 0 K.

\textbf{Classification:} Design gap. Silent omission. Must be flagged at minimum.
\end{mdframed}

\subsection{E-7 \textemdash{} Bead Orientation Duplicated Across Layers}
\begin{mdframed}[backgroundcolor=amber!8]
\textbf{Location:} \texttt{coarse\_grain/core/bead.hpp} (\texttt{Bead::orientation}),
\texttt{v5/bead\_transport.hpp} (\texttt{BeadDynState::n\_hat}).

\textbf{Problem:} Orientation is stored in two separate structs with no
defined synchronisation. \texttt{Bead::orientation} is the inertia-frame
static axis. \texttt{BeadDynState::n\_hat} is the transport-layer orientation
axis, initialised from the scene bead but thereafter independently tracked.
After any dynamic step, the two diverge with no update protocol.

\textbf{Classification:} Medium. Two truths = future bug shrine.
Define one owner.
\end{mdframed}

\subsection{E-8 \textemdash{} Gradient Verifier Print Bug}
\begin{mdframed}[backgroundcolor=gray!8]
\textbf{Location:} \texttt{src/sim/optimizer.hpp}, \texttt{verify\_gradients()}, line 363.

\textbf{Problem:} The numeric gradient diagnostic print re-calls
\texttt{model.evaluate\_energy(coords)} and prints the total energy,
not the numeric gradient component at the max-error index.
Only affects debug output when \texttt{check\_gradients = true}.

\textbf{Classification:} Low. Diagnostics only.
\end{mdframed}

\subsection{E-9 \textemdash{} TBP Axial/Equatorial Classification is Frame-Dependent}
\begin{mdframed}[backgroundcolor=gray!8]
\textbf{Location:} \texttt{src/pot/vsepr\_geometry.hpp}, \texttt{TRIGONAL\_BIPYRAMIDAL} case.

\textbf{Problem:} Classification of axial vs equatorial positions uses
$|dz_i|/r_i > 0.8$ --- a z-axis alignment test. This breaks for any
trigonal bipyramidal molecule that is not pre-aligned to the z-axis.
Since the FIRE optimizer does not enforce frame alignment, the
equatorial/axial assignment is not invariant under rotation.

\textbf{Classification:} High correctness risk. Fix requires
frame-independent classification (e.g., via inertia tensor decomposition).
\end{mdframed}

\subsection{E-10 \textemdash{} VSEPR Domain Gradient Jacobian Approximation}
\begin{mdframed}[backgroundcolor=gray!8]
\textbf{Location:} \texttt{src/pot/energy\_vsepr.hpp}, lines 134--141.

\textbf{Problem:} The gradient propagation from domain-repulsion through
bond-pair directions approximates the full Jacobian
$\partial\hat{\mathbf{d}}/\partial\mathbf{r} = (I - \hat{\mathbf{d}}\hat{\mathbf{d}}^\top)/|\mathbf{r}|$
by using only the $1/r$ magnitude scaling. The transverse projection tensor
$(I - \hat{\mathbf{d}}\hat{\mathbf{d}}^\top)$ is not accumulated.
This is accurate near equilibrium where bond directions are slowly varying,
but degrades when bonds are significantly stretched.

\textbf{Classification:} Medium. Document as known approximation.
\end{mdframed}

% ==============================================================================
\section{Immediate Action Table}
\label{sec:action_table}
% ==============================================================================

\begin{center}
\begin{tabular}{llll}
\toprule
ID & Issue & Priority & Action \\
\midrule
E-1 & Two Vec3 types          & \textcolor{red1}{\textbf{P1 Blocker}} & Unify or bridge before refactor \\
E-2 & evolve() stub           & \textcolor{red1}{\textbf{P1 Blocker}} & Wire minimal empirical loop \\
E-3 & sane() ignores F[]      & \textcolor{amber}{\textbf{P2}}        & Add F.size() == N check \\
E-4 & T[N] undefined update   & \textcolor{amber}{\textbf{P2}}        & Mark diagnostic-only \\
E-5 & Bead mass = 1.0         & \textcolor{amber}{\textbf{P2}}        & Add conversion bridge metadata \\
E-6 & StateC metadata silent  & \textcolor{amber}{\textbf{P2}}        & Add completeness flags \\
E-7 & Orientation duplicated  & \textcolor{amber}{\textbf{P2}}        & Define one owner, freeze other \\
E-9 & TBP frame-dependent     & \textcolor{amber}{\textbf{P2}}        & Inertia-tensor classification \\
E-8 & Gradient verifier print & \textcolor{gray1}{\textbf{P3}}        & Fix numeric print path \\
E-10 & VSEPR Jacobian approx  & \textcolor{gray1}{\textbf{P3}}        & Document; fix post-stability \\
\bottomrule
\end{tabular}
\end{center}

% ==============================================================================
\section{Pre-Refactor Checklist}
% ==============================================================================

Complete this list before touching the architecture:

\begin{itemize}
  \item[$\square$] Freeze version tag \texttt{v4.0.4.05} with audit notes.
  \item[$\square$] Choose canonical Vec3 (\texttt{vsepr::Vec3} recommended as base).
  \item[$\square$] Add \texttt{atomistic::Vec3 = vsepr::Vec3} or explicit bridge.
  \item[$\square$] Wire minimal \texttt{evolve()} loop: validate $\to$ clear F $\to$ eval F $\to$ VV step $\to$ update ledger $\to$ quality check $\to$ log events.
  \item[$\square$] Add \texttt{F.size() == N} check in \texttt{sane()}.
  \item[$\square$] Add \texttt{(T.empty() || T.size() == N)} check in \texttt{sane()}.
  \item[$\square$] Document \texttt{T[N]} as diagnostic-only (no update rule yet).
  \item[$\square$] Document \texttt{A[N]} as intentional derived quantity ($\mathbf{a} = \mathbf{F}/m$, not stored).
  \item[$\square$] Add bead mass conversion bridge: \texttt{BeadDynState::mass\_amu} field + scaling.
  \item[$\square$] Define orientation ownership: \texttt{Bead} = static/identity, \texttt{BeadDynState} = dynamic.
  \item[$\square$] Add metadata completeness flags to \texttt{StateC::from\_state()}.
  \item[$\square$] Fix gradient verifier print (E-8).
  \item[$\square$] Version tags: \texttt{v4.0.4.05} (kernel audit), \texttt{v5.0.0-beta6} (UFX migration start).
\end{itemize}

% ==============================================================================
\section{Version Labels}
% ==============================================================================

\begin{center}
\begin{tabular}{lll}
\toprule
Track & Tag & Description \\
\midrule
v4 patch & \texttt{v4.0.4.04} & Current (post last bump) \\
v4 patch & \texttt{v4.0.4.05} & Kernel audit and UFX migration preparation \\
v5 beta  & \texttt{v5.0.0-beta6} & Current (post last bump) \\
v5 beta  & \texttt{v5.0.0-beta7} & Post-repair: P1 blockers resolved, P2 bridges added \\
\bottomrule
\end{tabular}
\end{center}

\bigskip
\noindent\textbf{Decision rule:} beta6 is the audit snapshot.
beta7 opens only after E-1 and E-2 are resolved and the checklist is complete.

\end{document}
